\subsection{Attaching finite-cutoff paths on the selected stationary route}
\label{sec:path-attachment-producer}

This section records the selected stationary route from its horizontal
annulus estimate through finite-cutoff path attachment. It separates three
different jobs: producing the annulus estimate, using it in a one-step path
construction, and assembling the resulting transitions. Only the last job
is a finite induction. The quantitative inputs to the one-step construction
are stated at their actual interface and are not inferred from that
induction.

\subsubsection{Bounds for the one-step stationary path}

The background quadratic data $B$ (called a \emph{bed} below) comprise
the curvature metric and its root coefficient array. The parameter list
$p$ (the \emph{clip}) specifies the quadratic test action, including its
direct coefficient $p.\mathrm{direct}$. A path is a collection of row, loop, cubic, and quartic coefficients. For a
parameter list $p$, strip width $\tau$, coefficient bound $A$, and loop bound $\Lambda$, the
predicate $R.\mathrm{Bounds}(p,\tau,A,\Lambda)$ is the path budget:
each exponentially weighted row family has sum at most $A$; the loop
atom-budget family has sum at most $\Lambda$; and the cubic and quartic
atom-budget families each have sum at most $A$. In addition, every atom is
invariant, the row symbols satisfy the metric identity, row-symbol terms are
summable, and the loop coefficient and its first moment have summable terms
and sum to zero. The attachment predicate $\mathrm{Attached}(B,R)$ says
that the path coefficients at each plaquette and cover offset sum to the
bed's root entry. In particular, these are simultaneous global summability
conditions, not merely pointwise bounds.

The precise annulus estimate, denoted $\mathrm{Ann}(G,\sigma,
\varepsilon,C,Z,\tau,A)$, has the following quantifiers. It supplies an
integer $L_{\min}$ and $D_0\geq0$. For every odd side $L\geq L_{\min}$,
every side sequence $m_i\to\infty$, every bed $B$, clip $p$, path $R$, and
$\Lambda$, impose the hypotheses of Proposition~\ref{prop:ex-annulus-estimate}:
\[
 R.\mathrm{Bounds}(p,\tau,A,\Lambda),\quad
 \mathrm{Attached}(B,R),\quad B.\mathrm{metric.strip}(\sigma,C),\quad
 p.\mathrm{direct}\leq Z,
\]
along with $\mathrm{Floor}(B,p,n,\varepsilon,C)$ for every $n$ and
$\mathrm{Match}(B,m_i,G)$ for every $i$. Then there are one Fourier kernel $Y$,
nonnegative strip constants $v,E$, and a single error sequence $e_i\to0$
such that
\[
 \mathrm{Strip}(v,E,Y_{ij})\quad\hbox{for every }i,j,\qquad
 \mathrm{Rec}\!\left(Y,\mathrm{Pole}(p,R,L^{-1}),
          m_i^4e_i\right)\quad\hbox{for every }i,
\]
and $|\alpha^{\mathrm{ann}}_j(Y)+\lambda_G\log L|\leq D_0$ for every color $j$.
Here $\mathrm{Strip}$ means a bounded periodic holomorphic extension
with width $v$ and bound $E$, and $\mathrm{Match}$ denotes the
finite-volume quadratic matching condition for $B$. The pole Hessian
$\mathrm{Pole}$ and its finite-mesh recording condition $\mathrm{Rec}$
are specified in Subsection~\ref{sec:ex-annulus-recording}.
The latter is a finite-mesh identity up to the stated error; no passage to a limit in $i$ is part of this assertion.

Proposition~\ref{prop:ex-annulus-estimate} supplies this estimate on the
same fine bed, clip and path. Its Wilson reference and its Wilson-subtracted
correction are added before any finite-period recording is compared. The
fee bounds give one common response strip; the recording error is then the
exponentially small torus-mean error. The face estimates of
Subsection~\ref{sec:ex-annulus-fleet-faces} give the two radial inequalities
\[
 |g(x)|\leq\frac{A_1}{(\|x\|+\mu)^4},\qquad
 |g(x)-h_G(x)|\leq\frac{A_2}{\|x\|^3}
                   +\frac{A_2\mu^2}{\|x\|^6}
 \quad(\mu\leq\|x\|\leq1).
\]
Their constants precede the blocking side. The following elementary
integration is the last step giving its logarithmic drift constant.

\begin{lemma}[Radial logarithmic coefficient bound]
\label{lem:path-radial-log}
Assume $0<\mu\leq1$, $A_1,A_2\geq0$, and let $E_4\subset\mathbb R^4$
be a measurable finite-volume region containing the unit ball. Let $g$ be
real measurable and integrable on $E_4$, and let $h_G$ be integrable on the
annulus $I_\mu=\{x:\mu<\|x\|\leq1\}$. Assume the radial estimates
\[
 |g(x)|\leq\frac{A_1}{(\|x\|+\mu)^4}\quad(x\in E_4),\qquad
 |g(x)-h_G(x)|\leq\frac{A_2}{\|x\|^3}
                   +\frac{A_2\mu^2}{\|x\|^6}\quad(x\in I_\mu).
\]
Assume also the coefficient and pole identities
\[
 \alpha^{\mathrm{ann}}_j(F)=c_{\rm wt}\int_{E_4}g(x)\,dx,\qquad
 c_{\rm wt}\int_{I_\mu}h_G(x)\,dx=-\lambda_G\log\mu^{-1},
 \quad c_{\rm wt}>0.
\]
For the productive picture $F$, set
\[
 C_{\rm rad}(A_1,A_2)=A_1\left(\frac{\pi^2}{2}
       +\operatorname{vol}(E_4)\right)+3\pi^2 A_2.
\]
Then
\[
 |\alpha^{\mathrm{ann}}_j(F)+\lambda_G\log\mu^{-1}|
       \leq c_{\rm wt} C_{\rm rad}(A_1,A_2).
\]
\end{lemma}
\begin{proof}
Write $E_4$ for the finite-volume four-dimensional integration region,
$B_1=\{x:\|x\|\leq1\}$, $B_\mu=\{x:\|x\|\leq\mu\}$,
$I=\{x:\mu<\|x\|\leq1\}$, and $Q=E_4\setminus B_1$.
Polar coordinates give $\operatorname{vol}(B_r)=\pi^2r^4/2$.
On $B_\mu$, the first radial estimate implies
$|g(x)|\leq A_1\mu^{-4}$, so
$|\int_{B_\mu}g|\leq A_1\pi^2/2$. On $Q$, one has $\|x\|\geq1$,
and hence $|g(x)|\leq A_1$; therefore
$|\int_Qg|\leq A_1\operatorname{vol}(E_4)$.

On $I$, integrate the second estimate in radial coordinates. The two
radial integrals are
\[
 \int_I\|x\|^{-3}\,dx=2\pi^2(1-\mu),\qquad
 \int_I\|x\|^{-6}\,dx=2\pi^2\frac{\mu^{-2}-1}{2}.
\]
Thus
\[
 \int_I|g-h_G|\leq 2\pi^2A_2(1-\mu)
       +\pi^2 A_2\mu^2(\mu^{-2}-1)\leq3\pi^2A_2.
\]
The pole integral is exactly the logarithmic term:
$c_{\rm wt}\int_I h_G=-\lambda_G\log\mu^{-1}$, by the assumed pole identity. The Fourier formula for $\alpha^{\mathrm{ann}}_j$ is
$c_{\rm wt}\int_{E_4}g$. Splitting the integral over
$B_\mu\cup I\cup Q$ and using the three absolute bounds proves the
claim. The radial integrals were evaluated above; the pole identity is a separate hypothesis. No uniformity beyond the displayed radial estimates is used.
\end{proof}

\subsubsection{The stationary chart receipts}
\label{subsec:path-stationary-chart}

The chart data consumed by the path interpolation is produced at each fixed
finite volume. It is useful to give this construction separately: the
positive Hessian gap, critical chart, and installed-matrix approximation
enter the later path theorem as supplied data, but their source is an
implicit-function argument, not an additional limit assumption.

Throughout this subsection a kernel choice with penalty
$K_{\rm pen}>0$, an odd blocking side $L$ and a phase $j$ are fixed, and
$R_{\rm macro}=(j+3)\ell+4\ell$ denotes the macro radius, where
$\ell=L.\mathrm{val}$. A one-step datum consists of a coarse side $c$
and a fine side $f=c\ell$ together with the blocking maps of the kernel.
It is called larger when its patch satisfies the branch compatibility
hypothesis of the kernel and
\[
 100R_{\rm macro}\le f,\qquad 1<c .
\]
All statements below are for a larger one-step datum, an allowed forest
$F$ of its blocking ratio, and a local choice at the parent of each fine
edge. Configurations are taken in the forest slice, and forest directions
are extended by zero on the tree edges.

Start with the planted comparison for a coarse site $i$ and a fine edge
$e=(x,a)$, where $i$ is the parent of $e$. Write $\ell=L.\mathrm{val}$ as above;
let $p_i$ be the planted root frame and $P_i$ the local predictor. The
planting estimates give nonnegative constants $C_0,D_0$ and, with
\[
 \alpha_0=2\max(1,\mathrm{predictorCut}^{-1})(2C_0+D_0),\quad
 u_0=(8\ell+1)D_0\ell,\quad
 v_0=C_0\ell^2+\alpha_0\ell^{-2}+C_0\ell^2,
\]
the finite-volume edge estimate
\[
 \operatorname{chord}\!\left(U_e(\operatorname{replica}V)_e^{-1}\right)
 \leq (u_0+v_0)\,\operatorname{port}_i(U)
       +\mathbf 1_{\{i\ne i+\mathrm{parentStep}(e)\}}
          \operatorname{wire}(P_i,V,e).
\]
Here the first term is obtained by following the forest root walk: its
length is at most the patch radius, and the planted bound charges each
walk edge by $D_0\ell$ times the port. If the edge stays in its parent cell,
the lifted path closes at the root, so this is the entire estimate. If it
crosses to a neighboring parent, the endpoint root-to-frame chords each
cost at most $C_0\ell^2\operatorname{port}$, while the predicted edge has
chord at most $\alpha_0\ell^{-2}\operatorname{port}$. The triangle inequality
adds these three terms, and the residual is precisely the displayed wire
term. This is the two-case proof of the planted comparison; the tree walk
and crossing endpoint costs are not hidden in a generic locality claim.

For a fixed step, summing the squared port terms over coarse sites uses
Cauchy--Schwarz twice: first over the box offsets and plaquette types,
then over the box offsets in the site sum. Injectivity of the map from a
coarse site/plaquette pair to its embedded fine plaquette prevents extra
multiplicity. The resulting factor is exactly the product of the box
cardinality, the plaquette count, and the site count appearing in the
curl factor. The probe wire is summed over coarse bonds and paid
by the second factor, the probe factor. Combining the two with the floor
bound yields
\[
 \sum_e\|v_e\|^2\leq \beta_\ell\,\mathcal D(v)
                      +\gamma_\ell\,\mathcal P(v),
\]
where the coefficients are
\[
 \beta_\ell=(2C^2)|\operatorname{Site}(L)|\,|\operatorname{Axis}|
       (|X|\,|\operatorname{PlaqAxes}|)\,|X|,
 \qquad
 \gamma_\ell=2|\operatorname{Site}(L)|,
\]
with $X$ the fixed patch-offset set and $C$ the edge-comparison constant.
Thus $\ell$ is the blocking side cardinality, not the coarse torus side.
Here $v$ is a forest direction, extended by zero on the tree edges,
$\mathcal D$ is the squared physical color derivative and
$\mathcal P$ the squared predictor penalty. Here is the Hessian identity used in this comparison. Let $h$ be a
coarse variation and $v$ a forest-coordinate variation. Write
$e(h,v)$ for the fine variation obtained by differentiating the forest
lift at zero, and write $Q$ for the linearized block predictor. The
linearized predictor residual is
$c(h,v)=h-Qe(h,v)$. For two pairs $(h,v),(h',v')$, the tested action
has second derivative
\[
 D^2\mathcal T(0)[(h,v),(h',v')]
 =\sum_f\langle d_1e(h,v),M_nd_1e(h',v')\rangle_f
   +\mathrm{penalty}\,\langle c(h,v),c(h',v')\rangle.
\]
Indeed the Hessian of the fine main action is the first displayed
bilinear form by Lemma~\ref{lem:ex-soft-quadratic-hessian}, while its first
derivative vanishes, so the second derivative of the forest lift does
not contribute. The predictor penalty vanishes together with its first
derivative at the identity; its quadratic Taylor term is one half of
$\mathrm{penalty}$ times the squared norm of the predictor residual.
The second-derivative chain rule therefore gives the second term and no
higher derivative of the predictor. Taking $h=h'=0$ gives the forest
Hessian used above, whose first term is bounded below by
$a\mathcal D(v)$ by the all-period floor. Assume $a>0$, $\mathrm{penalty}>0$, and
$\beta_\ell,\gamma_\ell\geq0$. Set
\[
 \eta=\frac{\min(a,\mathrm{penalty})}{1+\beta_\ell+\gamma_\ell}.
\]
Then $\eta>0$, $\eta\beta_\ell\leq a$ and
$\eta\gamma_\ell\leq\mathrm{penalty}$, and the preceding two inequalities
give the strict bound gap
\[
 \eta\sum_e\|v_e\|^2\leq
 D^2\mathcal T(0)[v,v].
\]
This is the quantitative Hessian lower bound used by the chart construction.

The remaining chart receipts follow in four steps. First, the exponential
coefficient budgets in the path bound make each uncut row and atom series
converge uniformly on a complex neighborhood of the physical field. Termwise
analyticity and summability therefore give smoothness of the finite-volume
main action. The chart compatibility hypothesis supplies the required analytic
composition along every smooth field chart. Second, in a forest chart the
lifted field agrees near the origin with the ordinary field coordinates;
the tested action is the flat action composed with this lift. Differentiating
this germ identity and using the finite-volume Hessian identity gives the
prismatic second variation and the comptroller bilinear form.

Third, fix one allowed choice of local test and write $F(x,z)$ for the
tested action in coarse coordinates $x$ and forest coordinates $z$. At
$(0,0)$, $F=0$ and $D F=0$. The bound gap implies that the vertical Hessian
$D_z^2F(0,0)$ is positive definite, hence invertible in the finite
dimensional forest space. The implicit-function theorem applied to
$D_zF(x,z)=0$ gives a smooth branch $z=\rho(x)$ with $\rho(0)=0$.
By the branch compatibility hypothesis, the tested actions for any two
local choices agree on a neighborhood of $(0,0)$ in the full variables
$(x,z)$, not merely on the graph of $\rho$. Their vertical derivatives
therefore agree near zero, so the same $\rho$ is critical for each
local choice.

Fourth, differentiate the stationarity equation at zero:
\[
 D^2F(0,0)\big[(x,D\rho(0)x),(0,v)\big]=0
 \quad\text{for every forest direction }v.
\]
The comptroller formula identifies this bilinear form as the sum of the
bed metric term and the penalty pairing of the predictor residual. The
installed tower matrix satisfies that same linear equation by the diagonal
two-scale trial. Uniqueness, from invertibility of the vertical Hessian,
therefore yields
$D\rho(0)x=\mathrm{relative}(W,x)$, the guessed-chart identity. Finally,
the chain rule at the stationary branch removes the term involving
$D_zF(0,0)$; the same comptroller formula and the diagonal tower identity
give the target bed metric as the Hessian of the reduced action. This is the
heat identity: for every local choice, the tested action along $\rho$
and one half of the kinetic quadratic form of the coarse bed metric
are both real-analytic at the origin and have the same second
derivative there. It is an equality, not an estimate. These arguments prove the fixed-period stationarity,
installed-tangent, and reduced-Hessian identities used by the transition.

The uniformity is as follows. The number $\eta$ displayed above depends
only on the kernel, the blocking side $L$, the phase $j$ and the floor
constant $a$: the factors $\beta_\ell,\gamma_\ell$ involve only $L$, the
patch-offset set $X$ of radius $R_{\rm macro}$ and the edge-comparison
constant $C$ of the planting estimate, which is itself chosen for
$(L,j)$ and the kernel. Hence the same $\eta$ is a gap for every larger
one-step datum with these $L,j$ and kernel, every coarse side, every
forest, and every bed, clip and path satisfying the hypotheses. In
particular it is uniform in the period. By contrast the neighborhoods
of the origin on which the implicit-function theorem gives the branch
$\rho$ may depend on the finite volume, and no uniformity is asserted
for them.

The exact interface for the one-step construction is the following. Fix a
kernel choice, a phase $J$, positive $\varepsilon,w,\delta$, and
nonnegative $D,Z,C,Z'$. The one-step construction proved in Proposition~\ref{prop:path-one-step-production} chooses a reel $R_1$ satisfying
$R_1.\mathrm{Bound}(s)$, nonnegative $\kappa$, and positive $s$. For every
$\tau>0$ and $A\geq0$, assume
$\mathrm{Ann}(G,w,\varepsilon,D,Z,\tau,A)$. It then chooses
$d\geq0$, $\gamma>0$, a period threshold $u$, and functions $\beta(L)$,
$\Lambda_0(L)$. For every patch with phase $J$ and odd side at least $u$,
every pair of side sequences $m_i,n_i$ with $m_i\to\infty$, every path
and beds/towers, the following hypotheses are required:

\begin{itemize}
\item an armada whose every constituent step is a larger step, and the
  two bed hulls (metric and root strip bounds, floors, and sworn conditions);
\item an incoming path $R$ attached to the fine bed $B$, with
  $R.\mathrm{Bounds}(p,\tau,A,\Lambda)$ for the fine clip $p$, the
  diagonal two-scale trial of the fine bed $B$, the coarse bed $T$ and
  the tower, and the stated header bounds on the installed tower column;
\item for each forest chart $F_i$, the exact receipts
  \[
  \mathrm{HessLower}(F_i,p,R,\eta),\qquad
  \mathrm{Critical}(F_i,p,R,\rho_i),\qquad
  \mathrm{LinMatch}(F_i,W_i,\rho_i),
  \]
  and the heat identity for the coarse bed $T$;
\item the local-ratio chart and the specified direct-coefficient cutoffs
  for both clips. The annulus picture is produced inside this proof from
  $\mathrm{Ann}$; it is not an additional input.
\end{itemize}
The remaining scalar hypotheses are $1\leq\Lambda$, $0<\eta$, the phase
identity, the period threshold, and the two direct clip bounds. The armada
must have the selected kernel and all its patches must be branched. These
conditions are part of the one-step input, as are the hypotheses used to
define each forest chart.

The content of these hypotheses is precise. $\mathrm{HessLower}$ is the
positive lower bound
$\eta\sum_e\|v_e\|^2\leq D^2\mathcal T(0)[v,v]$ for the tested action
$\mathcal T$.
$\mathrm{Critical}$ asserts that the chart $\rho$ is smooth, vanishes at
zero, and is stationary for every tested action. $\mathrm{LinMatch}$ says
that the derivative of $t\mapsto\rho(tx)$ at zero is the chart's linear
prediction from the installed matrix. The preceding subsection proves these
receipts at each fixed finite volume by the planted sensitivity estimate,
the Hessian comparison, and the implicit-function theorem. The output
$\mathrm{ClassicalMatch}$ is instead a conclusion: for every Taylor order
$j\leq4$ and every local choice, the tested action along $\rho$ and the
main action of the new path have equal $j$th derivatives at the origin.

The one-step construction gives a path $Y$ and scalar $\alpha$ with
\[
 Y.\mathrm{row}=\mathrm{rowPath}(T),\qquad
 \mathrm{Attached}(T,Y),
\]
for the coarse bed $T$ and the coarse clip $q$; thus the output path is
attached to the coarse bed, while the incoming path was attached to the
fine bed. All four components of $Y.\mathrm{Bounds}(q,\eta',A',
\Lambda')$ obey the requested
row/loop/cubic/quartic budgets whenever
$0<\eta'\leq s$, $\kappa\leq A'$, $\Lambda_0(L)\leq\Lambda'$ and
$\Lambda/128\leq\Lambda'$. Under the same four conditions the loop,
cubic and quartic coefficient arrays of $Y$ are the images, under the
fixed encodings of the reel $R_1$, of elements $Q,P,W$ of its three
canonical coefficient subspaces for the coarse side sequence, with
\[
 \|Q\|\le\Lambda',\qquad\|P\|\le A',\qquad\|W\|\le A' .
\]
The reel consists of three coefficient Hilbert spaces with encoding
widths $\delta_2,\delta_3,\delta_4$, and $R_1.\mathrm{Bound}(s)$ means
that $s$ is strictly smaller than each of them. The canonical subspaces
are the orthogonal complements constructed in
Subsubsection~\ref{subsec:path-vow-representatives}. The construction also gives
\[
 |\alpha+\lambda_G\log L|+1\leq d,
\]
with the number $d$ chosen above.
It gives the order-four classical matching identity and, whenever
$\beta(L)\leq\Theta$ and $c'\leq\gamma$, the finite-cutoff bound with
amplitude $\Theta$ and decay rate $c'$. The latter is the following
statement at every index $i$, for the coarse side $c=m_i$ and the forest
$F_i$. The vertical Hessian determinant of the tested action is positive
along $\rho_i$ near the origin. The thermal function
\[
 \vartheta(x)=\mathcal L_R(x,\rho_i(x))
   +\tfrac12\log\det D_z^2\mathcal T(x,\rho_i(x))
   -\log\prod_b J(\rho_i(x)_b),
\]
formed with the loop term $\mathcal L_R$ of the incoming path, the
tested action $\mathcal T$ and the Jacobian $J$ of the local-ratio
chart, is real-analytic at the origin, and so is the outgoing
comparison function $\alpha\,\mathcal K_{2,Y}+\mathcal L_Y$ built from
the quadratic term and the loop term of $Y$ in the coarse clip. For
every local choice, every two coarse sites and every two polarization
amplitudes with complex entries of modulus at most one, the second
derivatives at the origin of these two functions on the corresponding
pair of real plane-wave directions differ by at most
\[
 |\operatorname{Site}(c)|\,\Theta\,e^{-c'c}.
\]
The order of the choices matters: the reel, $\kappa$ and $s$ are chosen
before $\tau,A$; the numbers $d,\gamma,u$ and the functions
$\beta,\Lambda_0$ after them; and the path $Y$ and the scalar $\alpha$
only after the incoming path, beds, tower and stationary data are
fixed. The proof constructs the classical
match by reducing the tested action to the new pole test and invoking
fourth-order exactness of the local response; the exponential estimate
comes from the funded Fourier ticket and the exponential rate comparison
below. The bounded path is obtained by adding the blank path to the produced
row and higher-degree coefficients, then applying the sum-budget estimate
to each degree separately.

The intermediate scalar comparison uses the determinant estimate of
Subsection~\ref{sec:ex-stationary-determinant-prices}, the contact cancellation
of Subsection~\ref{sec:ex-thermal-contact-cancellation}, and the alias
comparison of Subsection~\ref{sec:ex-quadratic-pole-transport}. These estimates
are applied to the same selected bridge. In particular the transverse
primal identity allows each nonprincipal nose to be replaced, in the
annulus pairing, by its horizontal part, which vanishes to second order.
This replacement, and not a second-order bound on the full nose, is the
one used below.

The scalar $\alpha$ is the local quadratic coefficient of the constructed
picture. Its normalization is compared with the annular coefficient
$\alpha^{\mathrm{ann}}_j(Y_\pi)$ by a finite response identity. The response identity
uses two determinant pictures. The first is the half-medal response minus
the reference determinant; the second is the half-rank-up response minus
the same reference determinant. The determinant comparison proves an estimate which
bounds the difference of their quadratic coefficients by
\[
 \left|\alpha^{\mathrm{ann}}_j(\text{medal}-\text{reference})
       -\alpha^{\mathrm{ann}}_j(\text{rank-up}-\text{reference})\right|
 \leq B_{\rm det},
\]
where $B_{\rm det}$ is the displayed finite reserve bound
$\mathrm{biteCost}(\eta/2,\mathrm{reserveCost}(\cdots))$. The rank-up
picture is recorded against the pulled-back pole test at the coarse meshes;
the annulus picture is recorded against the pole test itself at the cofinal
fine meshes, with error tending to zero. These are different finite-volume
tests, so their quadratic coefficients cannot be identified by comparing a
single test value. The exact comparison is an alias argument.

Let $L$ be the blocking ratio between the two periods, $\mu=L^{-1}$, $X(p)$ the analytic bridge
between the two period tests, $S$ the annulus picture, and $H$ the rank-up
picture. For a picture $P$, write $\mathcal R_P(p)$ for its matrix of
quadratic test coefficients. The finite-mesh recording identities and
period-stride approximation imply, for every real momentum $p$,
\[
 \mathcal R_H(p)=L^4\sum_{\ell\in\mathbb T_L}
  \mathcal S_S\!\left(\operatorname{alias}_\ell(p),
     N_X(-p,-\ell),N_X(p,\ell)\right),
\]
where $N_X$ are the bridge's nose maps and $\mathcal S_S$ is the annulus
matrix pairing. To prove this identity, choose actual mesh momenta $p_i\to p$
and stride representatives for each alias. The normalized recording error
for the annulus test tends to zero; the pulled-back rank-test error also
tends to zero after its period normalization. The bridge factorization
expresses each coarse test as the finite sum over its fine aliases. At a
fixed pair of aliases the recorded mixed coefficient converges to the
corresponding entry of $\mathcal S_S$. If the aliases do not pair as
$\ell=-k$, the two mesh centers remain distinct and the coefficient tends
to zero. The paired case gives the displayed entry. There are finitely many
aliases at this fixed period, so the limits pass through the sum. This is
the content of the finite-mesh transfer; it uses both recorded families,
not equality of their unsmeared tests.

The hypotheses used in this transfer are the ones present in the selected
rank/pole comparison: $X$ is analytic and has the required uniform strip;
the same matrix maps satisfy the exact coarse-to-fine premise; both metric
rows have their stated strip bounds; the fine bed has the required floor
and sworn conditions and the incoming path is attached to it; the rank
response is screened; the
cofinal recorded errors tend to zero; and the period meshes are related by
an injective replication map. These assumptions give the
normalized pulled-back recording and the noise-free alias identity. For the
Taylor step below, additionally use the exact second-bed identity
($N_X(0,0)=\mu I$), the transverse primal identity, and the annulus
picture's second-order profile. The nonprincipal alias estimate uses the
contractibility, strip bounds and horizontal-nose estimate proved in
Subsection~\ref{sec:ex-quadratic-pole-transport}; it does not follow from
the pole recordings alone.

The factor $L^4$ is essential. Here $L$ is the blocking ratio between the two periods, not the coarse physical period. Let $\ell=0$ denote the principal alias.
The second bed identity fixes the zero-momentum bridge nose at
$N_X(0,0)=\mu I$. The annulus picture has a second-order Taylor remainder,
so its principal pairing differs from $\mu^2\mathcal R_S(\mu p)$ by
$O(|p|^3)$. For a nonprincipal alias, the full nose need not vanish to second order. The
high-alias estimate in that comparison controls its horizontal and transverse components;
contractibility then replaces the full nose by its horizontal component
inside the pairing, without changing that pairing. The transverse primal
identity and the floor bound make this horizontal representative vanish to
second order. Only after this replacement does pairing with the annulus
quadratic profile give an $O(|p|^4)$ contribution. Summing the finitely many
aliases therefore gives
\[
 \mathcal R_H(p)-L^4\mu^2\mathcal R_S(\mu p)=O(|p|^3).
\]
Both sides are analytic near zero. Their Taylor coefficients through order
two must agree: along each ray, a nonzero coefficient of degree at most
two cannot be $O(|p|^3)$; polarization recovers the full quadratic form.
The second derivative of the rescaled principal term is
$L^4\mu^4D^2\mathcal R_S(0)$, and $L^4\mu^4=1$. Thus
\[
 \alpha^{\mathrm{ann}}_j(\text{rank-up}-\text{reference})=\alpha^{\mathrm{ann}}_j(Y_\pi).
\]
Thus $|\alpha-\alpha^{\mathrm{ann}}_j(Y_\pi)|\leq B_{\rm det}$. If the logarithmic
coefficient bound from the annulus is $D_0$, then
\[
 |\alpha+\lambda_G\log L|+1
 \leq B_{\rm det}+D_0+1.
\]
The path construction chooses $d=B_{\rm det}+D_0+1$ (or a larger bound).
This is the kinetic normalization step: it controls the finite-volume
quadratic coefficient by comparison with the continuum logarithmic
coefficient at the actual annulus scale. No unquantified limit in the
cutoff is used here.

The exact order-four match is proved from the fifth-order remainder ticket.
Let $f(x)$ be the tested action along the critical chart and let $g(x)$ be
the main term of the new path. The response construction gives smoothness
near zero and a ray estimate $f(tx)-g(tx)=O(|t|^5)$ for each direction $x$.
Taylor's theorem in the single real variable $t$ then forces the coefficients
of $t^j$ to agree for $j=0,1,2,3,4$. Those coefficients are the diagonal
values of the symmetric $j$-linear derivatives at zero. The polarization
identity for symmetric multilinear forms shows that equality on all
diagonals implies equality of the full derivatives. This proves
$\mathrm{ClassicalMatch}$ at every requested order; the input critical chart
and the bound gap are used to obtain the smooth response and the response
ticket, and are not deduced from the path identity.

For the path budget, the row component is the bed root array itself. Its
exponentially weighted row sum is bounded by
\[
 \mathcal B_{\rm row}(s',M)=|\mathrm{PlaqAxes}|\,M e^{4s'}
 \leq \mathcal B_{\rm row}(w,M),\qquad s'\leq w,
\]
because a row entry at offset $y$ is charged by the bed root strip and the
walk used to encode the slot has length $|y|_1+4$. The length convention
is the following. A slot consists of a walk $w$ on the cover and a
plaquette type; its length is $|w|+4$, the four counting the edges of
its plaquette. An atom with $r$ legs is an $r$-tuple of slots, and its
length is the sum of the lengths of its legs. Thus an atom whose legs
have walks of total length $|z|_1$ has length $|z|_1+4r$, and the
weight of an atom at width $\sigma$ is $e^{\sigma}$ to the power of
this length. The remaining three
components are the loop, cubic, and quartic atom families constructed by
the response step. Their weighted absolute budgets are separately bounded
by the three output bounds; invariance, zero loop jets, and exact metric
normalization are preserved by the construction. The sum-budget axiom
implies each family is summable, and the row-layout identity plus the floor
bound gives the metric clause. This is the explicit content behind the
claimed $Y.\mathrm{Bounds}$, not an appeal to pointwise coefficient control.

The finite-cutoff ticket is an absolute bound on the weighted remainder,
with original rate $k$ at least the requested rate $s$. If its unweighted
amplitude is bounded by $\Gamma$, and the final amplitude budget is
$\Lambda_1\geq1$, then for every cutoff $a\geq0$,
\[
 |wR|e^{-ka}\leq\Gamma e^{-ka}\leq\Gamma e^{-sa}
       \leq\Gamma\Lambda_1e^{-sa}.
\]
The proof uses only $s\leq k$ and nonnegative amplitudes. This is the
finite-cutoff exponential-rate comparison used in the one-step conclusion.

The path budget in the one-step construction is built in degree-specific steps. First, the blank
row path copies the target bed roots; if those roots have strip bound $M$,
its row budget is $\mathcal B_{\rm row}(w,M)=|\mathrm{PlaqAxes}|M e^{4w}$.
Second, the cubic and quartic coefficients are those of the classical
action restricted to its stationary graph, that is, of the tested action
along the critical chart; no determinant enters them. The cubic
coefficient is the cubic response of the tested action divided by $3!$.
The quartic coefficient is the quartic response, minus the three
contractions of two cubic responses through the inverse of the vertical
Hessian, one for each pairing of the four external legs, divided by
$4!$. This is the usual formula for the fourth derivative of a
stationary value.
If the cubic response bound is $A'$, the quartic response bound is $C'$,
and the inverse vertical Hessian is bounded at bound $D_*$, assume the hypotheses of Proposition~\ref{prop:ex-classical-stationary-pictures}, with its constants $C_3,C_4,F$ equal to $A',C',D_*$ respectively. That proposition gives, at its specified strip widths and with its Taylor normalization, the estimates
\[
 \frac{A'}{3!},\qquad
 \frac{C'+3A'D_*A'}{4!}.
\]
The factors $3$ and $3!$, $4!$ arise from the three pairings of four
derivatives and the Taylor factorials; each contraction is bounded by
$A'D_*A'$. These coefficients are recorded against the
third- and fourth-order Taylor coefficients of the tested action along
the critical chart, respectively. The determinant of the vertical
Hessian enters only the one-loop thermal function, hence the loop
correction below. Third, the loop correction
is obtained from the Fourier coefficient of the quadratic response after
subtracting its scalar quadratic term. The row path and these three atom
families are assembled into one path. Its metric identity follows from the
row-layout sum and the floor condition; loop normalization and invariance
are preserved coefficientwise. This is why the one-step construction has a
separate row bound and a separate total higher-degree bound, rather than
charging all coefficients to a single norm.

The final budget invoice in that estimate is elementary but essential. Let
$\gamma\geq0$, $V,\Lambda\geq0$, and suppose $\Lambda/128\leq\Lambda'$, the
terms arising before the annulus correction obey
\[
 2\bigl(C+(d+|y|)V\bigr)\gamma\leq\Lambda',
 \qquad |x+y|+1\leq d,
 \qquad B\gamma\leq\tfrac1{256}.
\]
Then
\begin{equation}
 \bigl(C+|x|V+B\Lambda\bigr)\gamma\leq\Lambda'.
 \label{eq:path-cap-invoice}
\end{equation}
Indeed $|x|\leq d+|y|$, by $x=(x+y)-y$, and
\[
 (C+|x|V+B\Lambda)\gamma
 \leq \bigl(C+(d+|y|)V\bigr)\gamma+B\gamma\Lambda
 \leq \Lambda'/2+\Lambda'/2.
\]
The final term uses $B\gamma\leq1/256$ and $\Lambda/128\leq\Lambda'$, so
$B\gamma\Lambda\leq\Lambda'/2$. This is
the only step where the auxiliary width of the horizontal picture is traded
against the outer row budget.

The exponential cutoff invoice is similarly explicit. If $\Lambda_1\geq1$
and $|wR|\leq\Gamma$, with $s\leq k$ and $a\geq0$, then
\begin{equation}
 wR e^{-ka}\leq \Gamma\Lambda_1 e^{-sa}.
 \label{eq:path-cutoff-invoice}
\end{equation}
The bound $|wR|\leq\Gamma$ gives $wR\leq\Gamma$, while
$e^{-ka}\leq e^{-sa}$; multiplying the latter by the nonnegative
$\Gamma\Lambda_1$ and using $\Lambda_1\geq1$ proves the claim. This is the
rate comparison used to preserve a finite-cutoff tail when the controlling
rate is weakened.

\subsubsection{Canonical representatives for all-period path jets}
\label{subsec:path-vow-representatives}

The path realization uses a closed-subspace projection to choose one
representative of prescribed jets at every tested period. Let $H$ be a real
Hilbert space, $S\subset H$ a closed subspace, and for each integer $n\geq0$
let $J_n:H\to E_n$ be a continuous linear map into a real Hilbert space
$E_n$. Define the joint invisible subspace
\[
 K=S\cap\bigcap_{n\geq0}\ker J_n,
 \qquad \mathcal R=S\cap K^\perp.
\]
Both are closed: each kernel is closed, and intersections of closed subspaces
are closed. Given $C\in S$, orthogonally project onto $K$ and put
$C_* = C-\Pi_K C$. Since $\Pi_K C\in K\subset S$, one has $C_*\in S$;
projection orthogonality gives $C_*\in K^\perp$, and hence $C_*\in\mathcal R$.
As $J_n(\Pi_K C)=0$, every prescribed jet is preserved,
$J_n(C_*)=J_n(C)$, and the Pythagorean identity gives
$\|C_*\|^2=\|C\|^2-\|\Pi_K C\|^2\leq\|C\|^2$.

This representative is unique. If $C_*,D_*\in\mathcal R$ have the same
jets for every $n$, then $C_*-D_*\in S$ and all its jets vanish, so it lies
in $K$. It also lies in $K^\perp$, whence
$\|C_*-D_*\|^2=0$ and $C_*=D_*$. Apply this with the admissible correction
series as $S$ and the derivatives at the tested periods as the maps $J_n$.
For this application assume the admissible subspace is closed; the joint invisible subspace is the common kernel used in the orthogonal-complement selection.
When the tested side sequence exhausts all physical periods, this proves
all-period jet injectivity on the chosen orthogonal complement. When it does
not, the conclusion is exactly for the specified test sequence. This is the
canonical realization argument: it uses the orthogonal complement of the
common jet kernel, not an assumed injectivity of the jet map on the original
space.

\subsubsection{Changing the forest without changing the stationary action}
\label{subsec:path-forest-independence}

\begin{lemma}[Forest-independent stationary germs]
\label{lem:path-forest-independence}
Fix one larger block step and its incoming path and clips. Let $F,F'$ be
two allowed forests, and let $\rho_F,\rho_{F'}$ be their critical branches
with positive vertical Hessians at zero. Their reduced classical actions
agree as germs of the coarse variables. Their thermal functions agree up
to a constant. In particular their derivatives of every positive order
agree on a sufficiently small common neighborhood of zero.
\end{lemma}
\begin{proof}
A configuration near the trivial configuration has a unique rooted gauge
transformation putting it into the slice of $F'$: start with the identity
at each root and solve successively along its tree edges. Each new gauge
value is a product of the already determined value and the intervening link
or its inverse. Thus the gauge transformation is analytic near the trivial
configuration, fixes the roots, and equals the identity there. Applying it
to a configuration in the $F$ slice defines an analytic change of vertical
coordinates $z'=\phi(x,z)$, with the coarse variable $x$ fixed. Repeating
the construction with $F$ gives its inverse. Shrinking the neighborhoods
if necessary, both transformations and their inverses are defined. Rooted
gauge invariance of the tested action gives
\[
       T_F(x,z)=T_{F'}(x,\phi(x,z)).
\]
The same holds for the inherited loop, since each loop atom is invariant.

The vertical chain rule shows that $\phi(x,\rho_F(x))$ is a critical
point for $T_{F'}$. It equals zero at $x=0$. The positive vertical Hessian
makes the implicit-function branch unique near zero; hence
$\phi(x,\rho_F(x))=\rho_{F'}(x)$ as germs. This proves the assertion
about the classical action and identifies the stationary representatives
before comparing determinants.

For completeness, the measure on either forest slice is the product Haar
measure on its free edges. This follows by integrating out the rooted
vertex gauges in order of tree height: at each stage one replaces a tree
link by a vertex gauge variable by a left or right multiplication, which
preserves probability Haar measure. The remaining free links are likewise
changed only by left and right multiplications. Fubini applies to this
finite product of probability measures. Applying this calculation first
to $F$ and then to $F'$ shows that the two slice measures are the same
measure on rooted gauge orbits. In local coordinates, if $j_F,j_{F'}$
are their positive density factors, the ordinary change-of-variables
formula therefore gives
\[
 c_Fj_F(z)=c_{F'}j_{F'}(\phi(x,z))
                     |\det D_z\phi(x,z)|,
\]
where the positive constants $c_F,c_{F'}$ are the coordinate normalization
constants, independent of $x,z$. Equality of continuous densities follows
from equality of their integrals over all compact coordinate sets. The
coordinate transformations may depend on $x$; this does not affect the
formula, which is on each fixed coarse fiber.

At the critical point, the second-derivative chain rule has no term from
$D_zT_{F'}$, and so
\[
 H_F=J^TH_{F'}J,\qquad J=D_z\phi(x,\rho_F(x)).
\]
Both Hessians are positive on a sufficiently small neighborhood. Taking
determinants, logarithms and the preceding density identity gives
\[
 \tfrac12\log\det H_F-\log j_F(\rho_F)
 =\tfrac12\log\det H_{F'}-\log j_{F'}(\rho_{F'})
                         +\log(c_F/c_{F'}).
\]
The inherited loop agrees on the two representatives. Adding it proves
the thermal assertion. In particular the second derivatives used in the
finite-cutoff comparison do not depend on the forest. No identification
of vectors in the two different forest coordinate spaces has been made.
\end{proof}

\subsubsection{Purchasing all components of the one-step path}
\label{subsec:path-one-step-production}

We now supply the one-step construction stated above. All widths in this
paragraph are widths in coarse momentum coordinates. The fine strip required
by the coefficient estimates is obtained by dividing them by $L$.

Fix the kernel, the phase, the two bed bounds, the positive floor, the
header bound and the two direct-coefficient ceilings in the one-step
interface. The selected bridge identities of
Subsection~\ref{sec:ex-selected-stationary-linear-data} and the weighted frame
and penalty estimates give one common positive coarse strip and fixed
bridge and penalty constants. The weighted inverse construction of
Subsection~\ref{sec:ex-weighted-stationary-inverses} narrows that strip and
bounds the needed inverse Hessians there. These choices depend only on
the listed fixed data, and precede the incoming path width $\tau$ and
bound $A$. The all-period bed floor supplies the real weighted coercivity in that
inverse construction. The forest gap of
Subsubsection~\ref{subsec:path-stationary-chart} separately gives each critical
branch; its possibly $L$-dependent numerical lower bound is not substituted
for the uniform weighted coercivity. Take the minimum of the finitely many
resulting widths and denote a further smaller positive width by $\sigma_0$.

For any subsequently fixed $\tau,A$, the residual coefficient bounds
of Proposition~\ref{prop:ex-greyson} have a positive fine width and finite
bounds. Enlarge the block-side floor until $\sigma_0/L$ lies in that
fine width. The classical vertex bounds are $a+L^{-1}\ell$ and
$b+L^{-1}r$, where $a,b$ were fixed before $\tau,A$, and $\ell,r$
are fixed after them. Enlarge the floor until $\ell/L,r/L\leq1$.
Proposition~\ref{prop:ex-classical-stationary-pictures} now gives the exact
cubic and quartic stationary pictures on one fixed coarse strip with the
bound $M_*$ displayed there. Both that strip and $M_*$ precede $\tau,A$.
The data used in that proposition are the actual bridge and the actual
inverse just described, so its three quartic contractions are the ones of
the critical branch of this step.

The selected kinetic response and differentiated matrix calculus of
Sections~\ref{sec:ex-selected-kinetic-response} and
\ref{sec:ex-stationary-determinant-prices} give the determinant picture
$f_L$ on a fixed smaller coarse strip. Its absolute bound $\Phi(L)$ may
contain the number $O(L^4)$ of fine variables; no bound independent of $L$
is asserted for it. Finite Fourier inversion identifies the normalized
finite-volume Hessian with the grid mean of this same kernel.
Lemma~\ref{lem:ex-pole-recording}, applied entrywise and then to the finite
color and direction sum, gives
\[
 \varepsilon_i\leq E(L)e^{-\gamma c_i},\qquad \gamma>0,
\]
where $c_i$ is the coarse side. The common coarse strip fixes $\gamma$
before $L$; the amplitude $E(L)$ may depend on $L$. The moving Haar
identity \eqref{eq:ex-thermal-determinant} identifies this picture with
the determinant and Haar part of the actual stationary thermal function.
It is therefore not a recording for a different coordinate action.

For the same incoming path, Proposition~\ref{prop:inherited-loop-picture}
gives an exactly recorded picture $g_L$ on a fixed coarse strip, of order
four and bound $C' L^{-2}\Lambda$. Its constant $C'$ may depend on $\tau$,
but not on $L$, $\Lambda$, the path or the periods. The extra fine-strip
condition in that proposition is met by increasing the same block-side
floor. The classical, determinant, inherited and outgoing bed pictures
are now all available on one smaller coarse strip independent of $\tau,A$
and $L$; only their thresholds and bounds have the dependence just stated.

Apply Proposition~\ref{prop:stationary-output-path} to these pictures and
the outgoing bed. Its coefficient-lift constants and widths consequently
precede $\tau,A$. The scalar is the quadratic coefficient $\alpha$ of
$f_L$. The core--rank comparison has an $L$-independent nuclear bound,
and the contact identity removes the actual--core quadratic difference.
Cauchy's estimate on the fixed strip bounds the resulting quadratic
coefficient difference by a constant $C_{\rm cmp}$ independent of $L$.
Proposition~\ref{prop:ex-pole-coefficient-transport} identifies the rank
coefficient with the fine annulus coefficient for this same bridge.
The annulus estimate of Proposition~\ref{prop:ex-annulus-estimate} thus gives
\[
 |\alpha+\lambda_G\log L|+1\leq
                d:=D_0+C_{\rm cmp}+1.
\]
The constant $d$ is chosen after $\tau,A$ and before $L$, as required.
In particular $|\alpha|\leq d+\lambda_G\log L$.

Write $B_H$ for the fixed outgoing quadratic-picture bound and $D$ for
the coefficient-lift bound multiplier, enlarging it also to pay for the
canonical coefficient realization described below. The loop budget before
canonical projection is at most
\[
 D\{\Phi(L)+(d+\lambda_G\log L)B_H+C'L^{-2}\Lambda\}.
\]
Increase the block-side floor once more so that $DC'L^{-2}\leq1/256$,
and put
\[
 \Lambda_0(L)=2D\{\Phi(L)+(d+\lambda_G\log L)B_H\}.
\]
If $\Lambda'\geq\Lambda_0(L)$ and $\Lambda'\geq\Lambda/128$, the
first two terms cost at most $\Lambda'/2$ and the last costs at most
$\Lambda/256\leq\Lambda'/2$. This proves the loop bound with precisely
the allowance in the one-step interface. The row, cubic and quartic
bounds are fixed before $\tau,A$, since the row bound and $M_*$ are.
Let $\kappa$ be their maximum, including the same realization multiplier.
The real-wave expansion in \eqref{eq:output-real-wave} changes the recording
amplitude only by $2(4\dim\mathfrak g)^2$; absorb it in $\beta(L)$.
The resulting classical match and finite-cutoff inequality are exactly
those of the one-step interface.

\subparagraph{Hilbert space of coefficients and canonical representatives.}
Here is an explicit Hilbert realization of the coefficient arrays, which
also fixes its choice before the incoming budgets. For each of the three
degrees use the canonical displacement arrays of the curvature lift,
together with their finitely many orientation transforms. A slot is encoded
by its displacement, its direction and color indices and a fixed coordinate
walk to that displacement. The number of such addresses of length at most
$h$ is bounded by a polynomial in $1+h$: the displacement lattice has fixed
finite dimension and the other index sets are finite. All coefficient
families produced by the lift are supported on these addresses. Choose
widths $0<s<\delta<\delta'<s_*$, where $s_*$ is a common width at which the
three lifted families have their proved absolute budgets. These widths
precede $\tau,A$.

For an address $a$ of length $|a|$, give its real coefficient space the
Hilbert norm
\[
 \|c\|_H=Z\left(\sum_a e^{2\delta|a|}|c_a|^2\right)^{1/2},
 \qquad Z=\left(\sum_a e^{-2(\delta-s)|a|}\right)^{1/2}<\infty.
\]
Finite-dimensional coefficient tensor norms are included in the address
set. The encoding is the coefficient array itself. Cauchy--Schwarz gives
\[
 \sum_a e^{s|a|}|c_a|\leq\|c\|_H,
 \qquad
 \|c\|_H\leq Z\sum_a e^{\delta'|a|}|c_a|.
\]
The second inequality uses $\delta<\delta'$ and $\ell^2\leq\ell^1$.
Thus a fixed enlargement of the previously obtained lift constants pays
for the Hilbert norm, whereas a Hilbert norm bound implies the desired
path budget at width $s$. The same construction with an intermediate
width in place of $s$ makes the encoding width strictly larger than the
final path width. This is the reel required in the interface.

The real, invariant and moment constraints define a closed real linear
subspace of this Hilbert space. Reality and invariance are intersections
of kernels of continuous coordinate functionals. The zero loop coefficient
and first moment are continuous sums, by Cauchy--Schwarz with the extra
exponential margin. At a fixed finite torus and fixed Taylor order, the
realization map is also continuous: derivatives of a fixed-order product
of transported plaquettes grow at most polynomially in its walk lengths,
and that polynomial is absorbed by the same exponential margin. Consequently
the invisible space for all requested jets at all the chosen periods is
closed. Project as in Subsubsection~\ref{subsec:path-vow-representatives}.
The projection preserves the real invariance and moment constraints and
every recorded jet, and decreases the Hilbert norm. For cubic atoms keep
both the third and fourth jets in this projection, for quartic atoms the
fourth jet, and for loops the second jet. Thus the same classical match
and the same thermal comparison survive simultaneously; changing the
cubic fourth jet would not have sufficed. All norm multipliers have been
included in $D$ and $\kappa$ above before choosing the last block-side
floor. This supplies the canonical representatives with the exact budgets
claimed in the one-step interface.

\subparagraph{The terminal path.}
At a terminal bed the tested quadratic kinetic form is the Wilson kinetic
form: its unamended diagonal is the identity, and the amendment identity
\eqref{eq:positive-amendment} gives $Gd_1=Hd_1=d_1$. Thus the amended
metric acts as the identity on exact curvatures, although it need not be
the identity on all two-forms. Take its root row as the row component of the path and denote its
quadratic action, including the direct clip term, by $J$. The floor and
sworn identities and Lemma~\ref{lem:ex-soft-quadratic-hessian} identify
$D^2J(0)$ with this kinetic form. Thus the Wilson action $W$ and $J$ have
the same value, first derivative and Hessian. The Wilson coefficient
pictures and the row-quadratic pictures supply exact common third- and
fourth-order pictures of $W-J$, with a width and bound depending only on
the terminal stage's root bound and direct-coefficient ceiling.
Apply the leading-removal construction of
Subsection~\ref{sec:stationary-output-construction} first at degree three
and then at degree four. In the second step subtract the fourth coefficient
of the cubic atoms as well. It gives cubic and quartic invariant arrays
with $W-J-S_3-S_4=O(|x|^5)$ at every tested period. Give the path these
arrays and zero loop coefficients. It is attached to the terminal bed,
its four path budgets hold below one fixed positive width and above one
fixed stage bound, and polarization gives equality of every derivative
through order four. No period-dependent selection or limiting argument
enters this terminal construction.

\begin{proposition}[Production of the one-step and terminal paths]
\label{prop:path-one-step-production}
Under the bed, tower, floor, attachment and header hypotheses of the
one-step interface, the annulus estimate supplies all its conclusions,
with the stated order of choices and the loop contraction $\Lambda/128$.
The terminal identity bed has the terminal base path required in
\textup{(B)} below. The classical match and thermal Hessian estimate hold
for every allowed forest.
\end{proposition}
\begin{proof}
The preceding construction verifies each input of
Proposition~\ref{prop:stationary-output-path} on the same beds, bridge,
critical branches and torus sequence. Its estimates, the explicit
split of the loop allowance, and the canonical projection give the
one-step path with all the required norms. The terminal construction
gives the base path. For each period choose any allowed forest to make
the output pictures. Lemma~\ref{lem:path-forest-independence} transfers
both matching assertions to every other forest, without changing the
output path or scalar. This proves the proposition.
\end{proof}

\subsubsection{The arbitrary-input factor-64 envelope}
\label{subsec:path-factor64-envelope}

The numerical factor used by the selected global caller is an allowance
amplification at the input to the transition supplier. We give both layers
of the finite reindexing, since the two maxima/minima use different pairs.

For each fixed output index $i$, the one-step theorem is first called at
both neighboring input indices $j=i$ and $j=i+1$. If the two calls return widths
$w_{i,i},w_{i+1,i}$ and row/cubic/quartic incoming-bound thresholds
$v_{i,i},v_{i+1,i}$, the common parameters are
\[
 s_i=\min(b_i,w_{i,i},w_{i+1,i}),\qquad
 A_i=\max(E_i,v_{i,i},v_{i+1,i}).
\]
Here $b_i,E_i$ are the width and bound supplied by the blank-path estimate.
For each of the two calls, $s_i$ is no larger than its width and its
bound threshold is no larger than $A_i$, so the same pair of finite-volume
estimates can be used at either adjacent input scale. The flux theorem is
universal in the admissible incoming path, bed, tower, armada, and chart;
therefore this uniformization does not choose a favorable incoming object.

Next fix the requested positive width sequence $\tau_j$ and nonnegative
incoming path-bound allowance $C_j$. For each fixed input index $j$, the
one-step calls indexed by the two possible output indices $i=j$ and
$i=j-1$ are unified by a second set of extrema: take the maximum of their
drift bounds, period thresholds, and side-dependent cutoff amplitudes, and
the minimum of their positive decay rates. At $j=0$, the predecessor is
interpreted as $0$, so the single available call is used twice. These choices
make one set of output data valid whether the selected transition uses the
same-scale receipt or the preceding-scale receipt. This is the second
neighbor reindexing in the construction; it is distinct from the
width/bound extrema at fixed output index above.

In the public application, the arbitrary allowance sequence $C_j$ is passed
as $64A'_j$, where $A'_j$ is any nonnegative sequence dominating the base
path-bound sequence. Thus the one-step conclusion is established for an
incoming path budget $64A'_j$; the factor 64 is part of the input allowance,
not an output estimate obtained from a smaller path. At the selected level,
$A'_j$ is the profile path bound, so the row, cubic, and quartic incoming
budgets are $64\,\mathrm{paths}_j$. The loop halo is separately defined with
budget $64\,\mathrm{loopCap}_j$. The selected schedule must satisfy the following
loop charges used by the transition. For the transition from level
$m+1$ to level $m$, let $j=J_{m+1}$ be the stage of the incoming level
and $L_j$ the block side selected at stage $j$; the stage of the output
level is $J_m\in\{j,j-1\}$. The side-dependent threshold $W_j$ of the
one-step construction at stage $j$, namely the function $\Lambda_0$ of
its interface, and the loop capacities satisfy
\[
 W_j(L_j)\leq\mathrm{loopCap}_m,\qquad
 \frac{64\,\mathrm{loopCap}_{m+1}}{128}
      =\frac{\mathrm{loopCap}_{m+1}}2
      \leq\mathrm{loopCap}_m,\qquad 1\le\mathrm{loopCap}_m .
\]
Thus the threshold is indexed by the incoming stage and the capacity by
the output level. These three inequalities are properties of the scale
schedule; they are inputs here. An
incoming path with ordinary profile bounds already satisfies the halo bounds
by componentwise monotonicity, since each budget is nonnegative and
$1\leq64$. That monotonicity only converts an already constructed path; the
arbitrary-input transition itself comes from the pairwise flux calls and the
two finite extrema operations. Its horizontal receipt and local response
bounds remain the explicit inputs described in the one-step section.

\subsubsection{From adjacent one-step paths to one finite history}

Let $a$ be the scale schedule, and write $J_m=a.\mathrm{clock}.J(m)$. Its
slow-step property gives
\[
 J_m\leq J_{m+1}\leq J_m+1.
\]
Suppose that for every pair $(j,i)$ with $j=i$ or $j=i+1$ the one-step
receipt is available, with width $w_{j,i}>0$, bound $A_{j,i}$, and path
$Y_{j,i}$. Define the common scale width and bound by
\[
 s_i=\min\{b_i,w_{i,i},w_{i+1,i}\},\qquad
 A_i=\max\{E_i,A_{i,i},A_{i+1,i}\}.
\]
Taking minima and maxima preserves positivity and nonnegativity. At either
adjacent index $j\in\{i,i+1\}$, one has $s_i\leq w_{j,i}$ and
$A_{j,i}\leq A_i$. Thus both adjacent-stage one-step constructions meet
the common width and bound requirements of output stage $i$. For each
actual transition the construction of the corresponding stage pair is
applied to its fixed incoming path and quadratic data. The extrema
unify the admissible parameters; they do not assert that the paths
returned by the two constructions coincide. The common bound is the
maximum of the two one-step bounds, and
the common cutoff rate is the minimum of the two one-step rates. This finite
neighbor maximum/minimum is the entire unification step in the global
construction; no limit in the scale index is taken here.

The next finite construction is where attached paths are assembled.
Branch rows and admissible paths are defined in
Definition~\ref{def:path-admissible}, and the transition rule on the actual
history is constructed in Proposition~\ref{prop:path-retained-plan}. Assume
that for each terminal index $N$ there is a branch row $R_N$, and for each
$m<N$ there is a transition rule which, from any admissible path at $m+1$,
returns a path at $m$, an amplitude $\alpha_m$, and a stationary coefficient
$\rho_m$, together with a bubble identity and the wing bound. Then for each
fixed $N$ one constructs rows $Y_m$ for $0\leq m\leq N$ by finite downward
induction. Set $Y_N=R_N$. Given $Y_{m+1}$, apply the transition rule to obtain
$Y_m$ and its amplitude/stationary data. Induction preserves admissibility
at the next stage because the previous step furnished exactly the attachment
and bound required as input. The resulting finite table has all adjacent
bubble identities and all wing bounds.

Doing this separately for each $N$ defines a single horizon of paths: its
$N$th column is the finite table just constructed. The terminal rows give the
branch-row property, the step outputs give the adjacent bubbles, and the
one-step row identities give the global row ties and path bounds. Hence one
obtains one finite-history path object satisfying all finite-cutoff path
conditions, with every cutoff retained on the same horizon. This argument is not a compactness extraction
and does not choose a different horizon at different cutoffs. It is the
finite descending construction furnished by the transition rule.

\subsubsection{The stationary plan on the retained horizon}
\label{subsec:path-retained-plan}

The finite descent above is conditional on a branch row and a transition
rule. We now state what these are on the actual history and how the
transition rule is obtained there. This is the statement consumed by the
later sections; its three analytic inputs are named explicitly.

Fix the scale schedule $a$. For a level $m$ write $J_m$ for its stage,
$p_m,\sigma_m,A_m$ for the clip, width and path allowance of the stage
profile, $\mathrm{loopCap}_m$ for the loop capacity, $L_m$ for the
block side from level $m+1$ to level $m$, and $n_{m,t}$ for the torus
side at level $m$ and period index $t$. A horizon $V$ assigns a bed
$V.H(N,m)$ to $m\le N$ and a tower $V.T(N,m)$ to $m<N$. Fix also a
local-ratio chart $O$ and post costs: positive numbers
$\eta_j,\Gamma_j,c_j$ for each stage $j$.

\begin{definition}[Admissible paths and branch rows]
\label{def:path-admissible}
A path $R$ is admissible at $(N,m)$, $m\le N$, if its row component is
the root row of the bed $V.H(N,m)$ in the walk encoding of the row-bound
estimate above, if $\mathrm{Attached}(V.H(N,m),R)$, and if
$R.\mathrm{Bounds}(p_m,\sigma_m,A_m,\mathrm{loopCap}_m)$. A branch row
at $N$ is a path $R$, admissible at $(N,N)$, such that at every period
the main action of $R$ in the clip $p_N$ and the Wilson action have
equal derivatives of every order $j\le4$ at the identity field, and
such that every loop coefficient of $R$ vanishes.
\end{definition}

\begin{definition}[Stationary plan]
\label{def:path-plan}
A stationary plan on $V$ consists of paths $R_{N,m}$ for $m\le N$,
numbers $\alpha_{N,m}$ for $m<N$, and maps $\rho_{N,m}^{t,F}$ for
$m<N$, every period index $t$ and every allowed forest $F$ of the
one-step datum from level $m+1$ to level $m$ at period $t$. It is
planned at the given post costs if the following hold for all $N$.
\begin{enumerate}
\item For $m\le N$ the path $R_{N,m}$ is attached to $V.H(N,m)$ and
 satisfies
 \[
 R_{N,m}.\mathrm{Bounds}(p_m,\sigma_m,A_m,\mathrm{loopCap}_m).
 \]
\item The terminal path $R_{N,N}$ has the Wilson derivatives of every
 order $j\le4$ at every period, and all its loop coefficients vanish.
\item For $m<N$ and all $t,F$, with $\rho=\rho_{N,m}^{t,F}$ and
 $j=J_{m+1}$:
 \[
 \mathrm{HessLower}(F,p_{m+1},R_{N,m+1},\eta_j),\quad
 \mathrm{Critical}(F,p_{m+1},R_{N,m+1},\rho),\quad
 \mathrm{LinMatch}(F,W,\rho),
 \]
 where $W$ is the installed matrix of the tower $V.T(N,m)$, and
 \[
 \mathrm{ClassicalMatch}(F,p_{m+1},p_m,R_{N,m+1},R_{N,m},\rho).
 \]
\item For $m<N$, $|\alpha_{N,m}+\lambda_G\log L_m|+1$ is at most the
 drift allowance of stage $J_{m+1}$.
\item For $m<N$ and all $t,F$, the finite-cutoff bound of the one-step
 interface holds for the incoming path $R_{N,m+1}$, the outgoing path
 $R_{N,m}$, the scalar $\alpha_{N,m}$, the chart $O$ and the map
 $\rho_{N,m}^{t,F}$, with the bound
 $|\operatorname{Site}(n_{m,t})|\,\Gamma_j\,\mathrm{loopCap}_{m+1}
 \,e^{-c_jn_{m,t}}$, $j=J_{m+1}$.
\end{enumerate}
\end{definition}

The paths of a plan are used twice: $R_{N,m+1}$ is the incoming path of
the step to level $m$ and the outgoing path of the step from level
$m+2$. It is one path with one choice of coefficient arrays.

The proposition below uses the following three constructions. Their
producers are given explicitly; the quadratic horizon is constructed in
its own chapter.
\begin{itemize}
\item[(A)] \emph{The annulus estimate.} $\mathrm{Ann}(G,\sigma,
 \varepsilon,C,Z,\tau,A)$ holds for all $\sigma,\varepsilon,\tau>0$
 and $C,Z,A\ge0$. Its producer is Proposition~\ref{prop:ex-annulus-estimate}.
\item[(B)] \emph{The one-step construction} with the interface stated
 in Subsubsection~\ref{subsec:path-stationary-chart}, for every stage pair
 $(j,i)$ with $j=i$ or $j=i+1$, and the terminal base path: for every
 stage, every bed with the root strip bound, the floors and sworn
 conditions of that stage and with diagonal equal to the identity at
 the tested periods, and every clip within the stage's direct bound,
 there is a path whose row is the root row of the bed, which is
 attached to the bed, has zero loop coefficients, satisfies the path
 bounds at every width below and every bound above fixed stage
 values, and has the Wilson derivatives of orders $j\le4$ at the
 tested periods. The one-step and terminal constructions are proved in
 Proposition~\ref{prop:path-one-step-production}; forest independence is
 Lemma~\ref{lem:path-forest-independence}.
\item[(C)] \emph{The quadratic horizons.} A free horizon $X$ and a
 retained horizon $V$ on the schedule $a$, and kernel bounds
 $(\alpha_j,\Omega_j,\eta^0_j)$, such that $V$ has the same beds as
 $X$ and, for $m<N$, the same installed column and the same transport
 column in its towers; and such that $X$ satisfies, for $m\le N$: the
 metric and root strip bounds of $X.H(N,m)$ at width $\alpha_{J_m}$
 and bound $\Omega_{J_m}$; the floor of $X.H(N,m)$ in the clip $p_m$
 with constant $\eta^0_{J_m}$ at every torus side; the sworn condition
 at every torus side; the identity $\operatorname{diag}X.H(N,N)=1$ at
 every side; and for $m<N$ the header bounds of the tower $X.T(N,m)$
 and the diagonal two-scale trial of $X.H(N,m+1)$, $X.H(N,m)$,
 $X.T(N,m)$ at every period.
\end{itemize}
The conditions on $X$ in (C) are the bed and flight receipts of the
quadratic construction of
Subsubsection~\ref{subsec:ex-lattice-horizon-producer}. The retained horizon
is formed from $X$ in Subsection~\ref{sec:common-funded-history}, where
$X$ and $V$ are called $M^0$ and $M^1$. Convergence along the cutoff
is not a consequence of (C): the conditions there do not determine
the remaining tower blocks. For the particular horizons constructed
jointly with their curvature roots it is proved in
Subsection~\ref{sec:free-horizon-convergence}; the present proposition
does not use it.

\begin{proposition}[Stationary plan on the retained horizon]
\label{prop:path-retained-plan}
Assume the gauge group is simple and \textup{(A)}, \textup{(B)} hold.
Fix a kernel choice, kernel bounds, and bounds $Z_j\ge0$ for the
direct clip coefficients of the stages. There are stage widths
$\bar\sigma_j>0$ and stage path bounds $\bar A_j\ge0$ such that, for
all stage widths $0<\sigma_j\le\bar\sigma_j$ and path allowances
$A_j\ge\bar A_j$, there are drift bounds $d_j\ge0$, rates
$\gamma_j>0$ and block-side floors $u_j$ with the following property.
Let the stage profiles have these widths and path allowances, drift
allowances at least $d_j$, block-side floors at least $u_j$ and direct
clip coefficients at most $Z_j$. Then, once the block sides
are chosen, there are positive gaps $\eta_j$ and volume thresholds, and
once the protection data are chosen there are entry delays, such that
for every paid schedule $a$ beyond these thresholds and delays, every
post costs with $\eta_j$ the given gaps, $\Gamma_j\ge1$ and
$c_j\le\gamma_j$, all horizons $X,V$ as in \textup{(C)} and every
local-ratio chart $O$, the following hold.
\begin{enumerate}
\item There is a stationary plan $(R_{N,m},\alpha_{N,m},\rho_{N,m}^{t,F})$
 on $V$, planned at the post costs in the sense of
 Definition~\ref{def:path-plan}. Every $R_{N,m}$ is admissible at
 $(N,m)$; in particular its row is the root row of $V.H(N,m)$. The
 finite-cutoff bound holds even with the factor $\mathrm{loopCap}_{m+1}$
 replaced by one.
\item For every level $m$ there is a reel, fixed before the cutoff and
 the period, whose three encoding widths are strictly larger than
 $\sigma_m$, such that for all $N>m$ the loop, cubic and quartic arrays
 of $R_{N,m}$ are the encodings of elements of its three canonical
 coefficient subspaces for the side sequence $t\mapsto n_{m,t}$, of
 norms at most $\mathrm{loopCap}_m$, $A_m$, $A_m$ respectively.
\item \emph{Extended domain.} Let $m<N$ and let $R$ be any path
 attached to $V.H(N,m+1)$ with
 \[
 R.\mathrm{Bounds}(p_{m+1},\sigma_{m+1},64A_{m+1},
 64\,\mathrm{loopCap}_{m+1}).
 \]
 Then there are a path $D$ admissible at
 $(N,m)$, a scalar and stationary maps satisfying items \textup{(3)--(5)}
 of Definition~\ref{def:path-plan} with $R$ as incoming and $D$ as
 outgoing path, and the arrays of $D$ lie in the canonical subspaces of
 the same reel with the same norm bounds as in the previous item.
\item For every level $m$ there are $v_m>\sigma_m$ and $C_m$ such that
 the root row of $V.H(N,m)$ has the strip bound at width $v_m$ and
 bound $C_m$ for all sufficiently large $N$.
\end{enumerate}
\end{proposition}
\begin{proof}
The stage widths and bounds are those of
Subsubsection~\ref{subsec:path-factor64-envelope}. Let $b_i>0$ and
$E_i\ge0$ be the width ceiling and bound floor of the terminal base
path of \textup{(B)} at stage $i$, and let $w_{j,i},v_{j,i}$ be the
widths and bound thresholds of the one-step constructions for the
stage pairs $j\in\{i,i+1\}$. As there, put
$s_i=\min(b_i,w_{i,i},w_{i+1,i})$ and
$A_i=\max(E_i,v_{i,i},v_{i+1,i})$, and take
$\bar\sigma_i=\min(s_i,\alpha_i/2)$ and $\bar A_i=A_i$. Thus
every admitted width is at most $b_i$ and every admitted allowance
at least $E_i$, as the terminal rows require; and they are at most
$w_{j,i}$ and at least $v_{j,i}$ respectively, as the transitions
require.
The one-step constructions are called with incoming allowance
$64A_j$, which gives $d_j,u_j$, rates $\gamma'_j>0$ and the functions
$\beta_j,\Lambda_{0,j}$ of the block side; the rates of the
statement are $\gamma_j=\gamma'_j/2$. The gaps
$\eta_j$ are those of Subsubsection~\ref{subsec:path-stationary-chart} for
the block side $L_j$, the phase $j$, the kernel and the floor constant
$\eta^0_j$. The volume threshold at stage $j$ is the maximum of
$100(j+7)$ and of a coarse side beyond which
$\beta_j(L_j)e^{-\gamma'_jn}\le e^{-\gamma_jn}$. The entry delays are those beyond which the three
loop-capacity inequalities of
Subsubsection~\ref{subsec:path-factor64-envelope} hold.

\emph{Terminal rows.} By \textup{(C)} the bed $V.H(N,N)=X.H(N,N)$ has
the root strip bound, floors, sworn condition and identity diagonal.
The base path of \textup{(B)} at stage $J_N$, taken at width $\sigma_N$,
bound $A_N$ and loop allowance $\mathrm{loopCap}_N\ge1$, is therefore a
branch row at $N$.

\emph{The transition rule.} Let $m<N$, put $j=J_{m+1}$ and $i=J_m$, so
$j\in\{i,i+1\}$, and let $R$ be attached to the fine bed
$B=V.H(N,m+1)$ with the factor-$64$ bounds of item (3). We verify the
hypotheses of the one-step construction of the stage pair $(j,i)$ for
the fine bed $B$, the coarse bed $T=V.H(N,m)$, the tower $V.T(N,m)$,
the fine and coarse clips $p_{m+1},p_m$, and the one-step data at all
period indices $t$, whose coarse and fine sides are $n_{m,t}$ and
$n_{m+1,t}$. The paid schedule gives that the block side is $L_j$, the
phase is $j$, the block side is at least $u_j$, and the sides
$n_{m,t}$ tend to infinity with $t$. The volume threshold makes every
one-step datum larger. The two bed hulls, namely the metric and root
strip bounds, all floors and all sworn conditions of $B$ and of $T$,
are \textup{(C)} at levels $m+1$ and $m$; here the beds of $V$ are
those of $X$. The header bounds and the diagonal two-scale trial at
every period are \textup{(C)} as well, since they involve only the
beds and the installed and transport columns, which $V$ and $X$ share.
For every $t$ and every forest $F$, the chart receipts of
Subsubsection~\ref{subsec:path-stationary-chart}, applied to $B$, $T$, the
tower and $R$, give $\mathrm{HessLower}(F,p_{m+1},R,\eta_j)$ and a map
$\rho^{t,F}$ with the critical, guessed and heat properties; their
hypotheses are the diagonal trial, the floor of $B$, the bounds and
attachment of $R$, the sworn condition of $B$ and the floor of $T$.
Choose one forest $F_t$ for each $t$. All hypotheses of the one-step
interface now hold with $\Lambda=64\,\mathrm{loopCap}_{m+1}\ge1$. It
returns a path $D$ and a scalar $\alpha$ with
$D.\mathrm{row}$ the root row of $T$ and $\mathrm{Attached}(T,D)$. By
the choice of the stage widths and bounds and the first two
loop-capacity inequalities, the conditions
$\sigma_m\le s_i$, $\kappa\le A_m$, $\Lambda_{0,j}(L_j)\le
\mathrm{loopCap}_m$ and $64\,\mathrm{loopCap}_{m+1}/128\le
\mathrm{loopCap}_m$ hold, so $D$ satisfies
$D.\mathrm{Bounds}(p_m,\sigma_m,A_m,\mathrm{loopCap}_m)$ and has the
canonical representatives of item (2). Hence $D$ is admissible at
$(N,m)$. The scalar satisfies the drift bound because $d_j$ is at most
the drift allowance.

The one-step construction gives the classical match and the
finite-cutoff bound only for the chosen forests $F_t$. They hold for
every forest by the forest-independence of the stationary data: for
two allowed forests $F,F'$ of a larger one-step datum, with critical
maps $\rho,\rho'$ and positive gaps for the same path, the tested
actions along $\rho$ and along $\rho'$ agree near the origin for every
local choice, and the two thermal functions differ by a constant near
the origin. The first statement transfers the classical match, whose
right side does not involve the forest. The second transfers the
finite-cutoff bound, which concerns second derivatives; positivity of
the determinant and analyticity of the thermal function hold for every
forest by the gap. This is precisely Lemma~\ref{lem:path-forest-independence}. Finally the amplitude
$\beta_j(L_j)$ and rate $\gamma'_j$ are replaced by $\Gamma_j\ge1$ and
$c_j\le\gamma_j$ using the volume threshold. This proves item (3), with the
factor one in the finite-cutoff bound.

\emph{Assembly.} An admissible path at $(N,m+1)$ satisfies the
factor-$64$ bounds, since all budgets are nonnegative and
$\mathrm{loopCap}_{m+1}\ge1$. The finite downward induction of the
preceding subsection, started from the branch row and using the
transition rule just constructed, gives for each $N$ the paths
$R_{N,m}$, scalars and maps of item (1) together with item (2). The
conditions of Definition~\ref{def:path-plan} are the recorded outputs
of the steps; the finite-cutoff bound is weakened by the factor
$\mathrm{loopCap}_{m+1}\ge1$. Item (4) holds with $v_m=\alpha_{J_m}$
and $C_m=\Omega_{J_m}$, by the root strip bound in \textup{(C)} and
$\sigma_m\le\alpha_{J_m}/2<\alpha_{J_m}$; it holds for every $N>m$.
\end{proof}

The proposition produces, on the retained horizon itself, the attached
paths with their bounds, and consumes the floors, sworn conditions,
metric strips and clip bounds of that same horizon. It is a statement
about global path attachment. It does not by itself give the local
two-bed statement used in the nonlinear step, which also requires the
roll identities and is constructed from the same horizon and paths in
Subsection~\ref{sec:ex-lattice-selected-course-binding}.

\subsubsection{Selected cutoff limits}
The ten retained bed and tower row families must converge on the lattice
cover along the actual terminal cutoff sequence. Subsection~\ref{sec:free-horizon-convergence}
proves this for the recursively selected free horizon, using its uniform
strip bounds, positive root choices and committed transition identities.
Its proof includes the high-alias and noise limits, the exterior lifts and
amendments, and recovery of cover coefficients from the unbounded selected
period sequence. Subsection~\ref{sec:selected-history-limits} then applies
these limits to the same stationary construction and shows that the
terminal Wilson normalization preserves every fixed-scale limit.

\subsubsection{Dependency order}

The annulus estimate is supplied by the common response and the radial
Fourier estimates. The one-step construction above then combines the
classical pictures, the determinant and inherited-loop pictures, and the
curvature coefficient lift; the terminal path is constructed independently
from the Wilson jets. Only afterwards does the retained-horizon proposition
perform its finite downward induction. The quadratic horizon and its
cutoff convergence remain the separate constructions of
Sections~\ref{subsec:ex-lattice-horizon-producer} and
\ref{sec:free-horizon-convergence}. No convergence of the nonlinear paths
is used to produce their finite-cutoff attachment or bounds. Their later
limits use the canonical all-period representatives and the separate
argument of Subsection~\ref{sec:selected-history-limits}.
